% Cote 119: offprints, reports, courses, CV and letters (typed and printed part): archive cote 119, site page 205 (PDF page 206).
% typescript. Draft compared: none.
% AI-transcribed by Claude Opus 5.5 (claude-opus-5-5, adaptive thinking, effort high), Message Batches API, batch r119_full (msgbatch_01DwZNm5Y7YsKBWyAN4sFs8w), 2026-11-13; not yet checked by eye.
\scanpage{119/205}
\mspage{- 10 -}

4°) \uline{Cohomologie cristalline} : quelques idées de départ sont esquissées
dans [18, Exp. IX], puis reprises et systématisées dans la thèse de P. Berthelot
[*], et dans le travail de P. Berthelot et L. Illusie sur les classes de Chern
cristallines[*].

e) \uline{Théorie du groupe fondamental} ([16, SGA 1, SGA 2, SGA 7, Exp. I et II],
[15 bis, n° 182], [19 bis]):

D'un point de vue algébrico-géométrique, tout était à faire, depuis
la définition du groupe fondamental d'une variété quelconque, en passant par
des propriétés ``de descente'' incluant des résultats assez formels du type de
van Kampen, jusqu'au calcul du groupe fondamental dans les premiers cas non
triviaux, comme celui d'une courbe algébrique privée de certains points; on
peut y adjoindre les théorèmes de génération et de présentation finie du groupe
fondamental d'une variété algébrique sur un corps algébriquement clos. Ce pro-%
gramme est accompli pour l'essentiel dans SGA 1, en utilisant à la fois les
résultats classiques sur le corps des complexes (établis par voie transcendante)
et une panoplie d'outils faits sur mesure (théorie de la descente, étude des
morphismes étales, théorème d'existence de faisceaux cohérents ...). Les
autres références contiennent des résultats plus spéciaux : théorèmes du type
de Lefschetz dans SGA 2, action des groupes de monodromie locale sur le groupe
fondamental d'une fibre dans SGA 7, Exp. I, calculs de certains groupes fonda-%
mentaux locaux dans [19 bis], via les groupes fondamentaux de certains schémas
formels. Tous ces résultats ont été utilisés couramment dans de nombreux tra-%
vaux, et en ont inspiré d'autres comme la thèse de Mme Michèle Raynaud [*].

f) \uline{Théorèmes de Lefschetz locaux et globaux} pour les groupes de Picard, le
groupe fondamental, la cohomologie étale, la cohomologie cohérente. Il s'agit
ici de la comparaison entre les invariants (cohomologiques ou homotopiques)
d'une variété algébrique et d'une section hyperplane. Les idées de départ sont
développées dans [16, SGA 2]. Cependant, pour des énoncés ``définitifs'', en ter-%
mes de conditions nécessaires et suffisantes, se reporter plutôt à la thèse
de Mme Michèle Raynaud [*] déjà citée.

g) \uline{Théorie des intersections et théorème de Riemann-Roch}:

La principale idée nouvelle, c'est qu'il y a presque identité entre
le groupe ``de Chow'' des classes de cycles sur une variété $X$ , et un certain
groupe de ``classes de faisceaux cohérents'' (tout au moins modulo torsion), à
savoir le groupe $K(X)$ (mentionné dans le par. 3). Dans un contexte modeste

\margin{205}
